\section{Phase estimates before key collisions}
\label{sec:phases}

The scalar recipes in \Cref{eq:recipe-scalars} couple the two
endpoints of each unit. We estimate the cross contractions that arise
when these scalar conditions are tested against profiles in the
effective spaces. These estimates will supply the compatible small
tables in \Cref{sec:preparation}.

Throughout this section different units are independent, while the two
endpoints of a unit may have arbitrary dependence. A marked unit is a
unit together with choices
\(\lambda_i\in C_i^{\mathrm{old}}\), \(i=1,2\), where
\(C_i^{\mathrm{old}}\) is its effective space after the first peeling.
The marks may depend on the whole unit. Put
\begin{equation}
 \Delta_O=U_1\lambda_1+U_2\lambda_2,
 \qquad u_{O,S}=\sum_{i\in S}u_i
 \quad(S\subseteq\{1,2\}).
 \label{eq:delta-phase}
\end{equation}
The effective-space rank bound and rank subadditivity give total
component rank at most \(2K_1\) for \(\Delta_O\). We use the slightly
larger budget \(r=2K_1+2\).

One case explains why the four bits
\[
 u_{B,1}(\Delta_A),\quad u_{B,2}(\Delta_A),\quad
 u_{A,1}(\Delta_B),\quad u_{A,2}(\Delta_B)
\]
are relevant. Suppose, for this calculation, that
\(p_i=T(\lambda_i,\cdot)+a\) on \(\cX\), where
\(p_i=u_{i'}U_i\), and that
\(c=a(\lambda_1)+a(\lambda_2)\) is common to the units.
Test the recipe using the actual small table. Its effective-space
equation and symmetry of \(T\) give, for an endpoint \(z\) of the
opposite unit,
\[
 (p_i+a)(v_{iz})
 =T(\lambda_i,v_{iz})
 =a(\lambda_i)+(u_zU_i)(\lambda_i).
\]
The coupled recipe equation is therefore
\[
 0=\sum_{i=1}^2(p_i+a)(v_{iz})
   =c+u_z(\Delta_A).
\]
The reciprocal equations require all four displayed bits to equal
\(c\). More general descriptions of the \(p_i\)'s lead to parities
of the same four bits; \Cref{sec:preparation,sec:status-span} identify
the required parities. Here we prove estimates for arbitrary marks,
without assuming the identities used in this illustrative case.

We use \(L_0,d_0,u_0,K\) from
\Cref{eq:law-budgets,eq:pin-budgets}; in particular,
\[
 L_0=\lceil100(D+10)\rceil,
 \qquad d_0=4rL_0,
 \qquad u_0=10(K_1+d_0+1).
\]
All fixed positive mass restrictions in the phase argument have lower
bounds depending only on \(D,g,K_1\) and fixed numerical tolerances.
Their marginal densities are consequently at most \(L\mu\), for an
\(L\) depending only on these early parameters, and their joint
densities are at most \(2^{(D+.01)N}\mu^2\) for sufficiently large
\(n\). Fixed multiplicative factors affect only the lower bound on
\(n\).

\begin{definition}
A componentwise cover is a collection of subspaces
\(T_e^\pm\subseteq V_e^\pm\). It covers a tensor
\(\xi=(\xi_e)_e\) if
\[
 \xi_e\in T_e^+\otimes V_e^-+V_e^+\otimes T_e^-
 \quad\text{for every }e.
\]
Its two mode dimensions are
\(\sum_e\dim T_e^+\) and \(\sum_e\dim T_e^-\); its total dimension
is their sum.
\end{definition}

\begin{theorem}[Phase alternative]
\label{thm:phase-alternative}
Suppose a marked-unit law is obtained by restricting an admissible
law in \Cref{thm:raw-supersaturation} to mass at least \(10^{-12}\),
and \(\Delta_O\ne0\) throughout. For all sufficiently large \(n\),
one can retain a sublaw of fixed positive relative mass, bounded below
using only the early parameters, such that one of the following holds
for independent draws \(A,B\):
\begin{enumerate}[label=\textup{(\roman*)}]
\item every nontrivial phase mean
\[
 \E(-1)^{u_{B,S}(\Delta_A)+u_{A,R}(\Delta_B)},
 \qquad(S,R)\ne(\varnothing,\varnothing),
\]
has absolute value less than \(.005\);
\item a componentwise cover of total dimension at most \(d_0\)
covers every \(\Delta_O\), and every displayed mean for which
\(|S|=1\) or \(|R|=1\) is \(2^{-\Omega(N)}\).
\end{enumerate}
The cover may depend on the law and on \(n\). The constants in its
dimension bound and in the relative mass bound do not depend on the
later selector or channel sizes.
\end{theorem}

In the first alternative, Fourier inversion gives every prescribed
four-bit pattern probability at least
\((1-15(.005))/16\). In the second, only the characters using an even
number of bits from each unit pair remain uncontrolled. The full-sum
kernel \(F(A,B)=u_{B,\{1,2\}}(\Delta_A)\) has rank \(O(N)\)
because \(\Delta_A\) belongs to the cover tensor space.
\Cref{lem:tiny-cover} uses this rank bound to control the remaining
phase conditions. When \(F(A,A)=0\), it gives an inverse-polynomial
lower bound for the event that all four bits equal a prescribed
\(c\). This smaller probability scale requires a separate final
overlap argument.

The proof first tests for a cover whose mode dimensions may grow
linearly with \(N\). If no such cover carries appreciable mass, the
following rank-growth lemma gives cancellation by a high-moment
estimate. If such a cover exists, its offset coordinates permit the
singleton estimates. A second concentration step either reduces that
cover to bounded dimension or gives cancellation for the remaining
characters as well.

\begin{lemma}[Rank growth without a cover]
\label{lem:no-cover-rank}
Let \(r\ge1\) be an integer, set \(m_r=100r+10\) and
\(q_r=2^{-m_r}\), and let \(s\) be a power of two with
\(s\ge2^{m_r}\).  Suppose \(T_1,\ldots,T_s\) are independent copies
of a random tensor of total component rank at most \(r\).
If every fixed componentwise cover with dimension at most \(rs\)
in each mode has probability less than \(p\), then
\begin{equation}
 \PP\left\{\rank\left(\sum_{j=1}^sT_j\right)<q_rs/4\right\}
 \le 4^s p^{q_rs/2}.
 \label{eq:no-cover-rank}
\end{equation}
Here rank means the sum of component ranks.
\end{lemma}

\begin{proof}
Regard all components as blocks in the direct sums of the two mode
spaces.  Their column and row spaces are direct sums of
componentwise spaces, so that every cover obtained in this proof is
componentwise.  We first make a deterministic observation about any
realized list of tensors.

Expose a subset \(E\) of the indices and quotient the two modes by
the column and row spans of the exposed tensors.  If \(b=s-|E|\)
indices remain, let \(C_j,R_j\) be their projected original column
and row spaces.  Define their span deficits by
\[
 \delta_C=\sum_{j\notin E}\dim C_j-\dim\sum_{j\notin E}C_j,\qquad
 \delta_R=\sum_{j\notin E}\dim R_j-\dim\sum_{j\notin E}R_j.
\]
The projected sum has rank at least
\begin{equation}
 \#\{j\notin E:\overline T_j\ne0\}-\delta_C-\delta_R.
 \label{eq:rank-deficits}
\end{equation}
Indeed, before adding the mode spaces, place the projected tensors
on the separate formal summands \(C_j\otimes R_j\).  Their block
diagonal sum has rank at least the displayed number of nonzero
tensors.  Applying the two addition maps can decrease rank by at
most their kernel dimensions, which are \(\delta_C,\delta_R\).

There is an exposed set leaving \(b=s/2^j\) indices for some
\(0\le j<m_r\) such that
\(\delta_C+\delta_R\le b/4\).  To prove this, order the realized
direction spaces by a uniform random permutation.  In either mode,
let \(a_t\) be the expected increase of their span dimension at the
\(t\)-th step.  Submodularity and exchangeability give
\[
 r\ge a_1\ge a_2\ge\cdots\ge a_s\ge0.
\]
For \(b_j=s/2^j\), set
\[
 f_j=a_{s-b_j+1},\qquad
 v_j=f_j-\frac1{b_j}\sum_{t=s-b_j+1}^s a_t.
\]
Each remaining index has expected projected dimension \(f_j\)
immediately after the first \(s-b_j\) exposures.  Consequently its
mode's expected deficit divided by \(b_j\) is \(v_j\).
The average increment on the first half of this tail is at least
\(f_{j+1}\).  It follows that
\[
 v_j\le f_j-f_{j+1}+\tfrac12v_{j+1}.
\]
Summing for \(0\le j<m_r\) yields
\[
 \sum_{j=0}^{m_r-1}v_j
 \le2(f_0-f_{m_r})-v_0+v_{m_r}\le2r.
\]
The two modes therefore have total expected deficit ratios at most
\(4r\) across these \(m_r\) scales.  Some scale has expected ratio
less than \(1/4\), and some permutation realizes a ratio at most
\(1/4\), as claimed.

If the original sum has rank less than \(q_rs/4\), its projection
at this exposed set has no larger rank.  By
\Cref{eq:rank-deficits}, at most
\(q_rs/4+b/4\le b/2\) remaining tensors have nonzero projections.
At least \(q_rs/2\) of them thus belong to the cover formed by the
exposed column and row spans.  Both cover dimensions are at most
\(rs\).

For fixed disjoint index sets \(E,J\), with
\(|J|=\lceil q_rs/2\rceil\), condition only on the tensors indexed
by \(E\).  The cover is then fixed, and the tensors indexed by \(J\)
are still independent.  The probability that all are covered is
at most \(p^{|J|}\).  There are at most \(2^s\) choices for each
index set.  Taking their union proves \Cref{eq:no-cover-rank}.
\end{proof}

\begin{proof}[Proof of \Cref{thm:phase-alternative}]
Write
\[
 c_s=10^{-12}(1+r)^{-4},\qquad q_0=2^{-100r-10},
 \qquad \theta=10^{-4},\qquad \Gamma=D+2r+10,
\]
and let \(s\) be the largest power of two not exceeding \(c_sN\).
For large \(n\), it satisfies the size requirement of
\Cref{lem:no-cover-rank} and \(s\ge c_sN/2\).
Choose \(p_0>0\), using only these early parameters, so small that
\begin{equation}
 \frac{q_0}{2}\log_2(1/p_0)
 \ge 2+\log_2(1/\theta)+\frac{2\Gamma}{c_s}.
 \label{eq:early-p0}
\end{equation}
When the channel size is chosen, require also
\begin{equation}
 \frac{hq_0}{4}
 \ge \log_2(1/\theta)+\frac{2\Gamma}{c_s}.
 \label{eq:early-h-moment}
\end{equation}
Further lower bounds for \(h\) arising below will again involve
only early parameters.

\medskip\noindent
\emph{No large cover.}
Suppose no fixed cover of dimension at most \(rs\) in each mode
has mass at least \(p_0\).  Fix a candidate value \(\xi\) of
\(\Delta_B\), a nonempty \(S\), and a set \(R\).  Define the
unit-side sign \(\psi(A)=(-1)^{u_{A,R}(\xi)}\).  Against completely
iid channel columns at the endpoints selected by \(S\), expansion
of the even \(s\)-th moment gives
\begin{align}
 &\E_{\mathrm{iid\ channels}}
 \left|\E_A\psi(A)(-1)^{u_{B,S}(\Delta_A)}\right|^s
 \notag\\
 &\hspace{25mm}\le
 \E_{A_1,\ldots,A_s}
       2^{-h\rank(\Delta_{A_1}+\cdots+\Delta_{A_s})}.
 \label{eq:phase-moment}
\end{align}
For one channel pair the expectation of its bilinear character
on a tensor \(\xi\) is \(2^{-\rank\xi}\).  Independence over
channels and components proves this assertion; selecting both
endpoints only squares the factor.  The signs \(\psi(A_j)\)
have absolute value one and can be removed for an upper bound.
By \Cref{lem:no-cover-rank}, the last expression is at most
\[
 2^{2s}p_0^{q_0s/2}+2^{-hq_0s/4}.
\]
The channel marginal comparison in \Cref{lem:frame-images}
multiplies this bound by a constant depending on \(h\) and the
number of components, but independent of \(n\).
Markov's inequality and
\Cref{eq:early-p0,eq:early-h-moment} bound the probability that
the inner mean has magnitude at least \(\theta\) by
\(C_h2^{-\Gamma N}\).

There are at most \(C_r2^{2rN}\) tensors of total component rank
at most \(r\): choose their component ranks, then factor each
matrix through that rank.  The component-rank choices and the
overcounting factors are bounded independently of \(n\).
Take a union over these candidate values and the bounded choices
of \(S,R\), and then pay the joint-law density
\(2^{(D+.01)N}\).  The exceptional \(B\)-mass is exponentially
small.  On all other \(B\)'s the conditional \(A\)-mean is at most
\(\theta\), even after inserting the actual value \(\Delta_B\).
Every required phase mean is therefore at most
\(\theta+2^{-\Omega(N)}<.005\).  Interchanging \(A,B\) deals
with \(S=\varnothing\).  This proves the first alternative
without a further restriction.

\medskip\noindent
\emph{A large cover and its offsets.}
Otherwise retain mass at least \(p_0\) in a fixed cover
\((T_e^+,T_e^-)\) of dimension at most \(rs\) in each mode.
Each individual primal frame span avoids these fixed subspaces
except on exponentially small mass.  Indeed its dimension is
\(\dim\cB\), its location is uniform under \(\mu\), and
\[
 \PP_\mu\{\im P_{ie}\cap T_e^+\ne0\}
 \le 2^{\dim\cB+\dim T_e^+-N+1};
\]
the same estimate holds for the minus mode.  Here
\(rs\le rc_sN\) and \(\dim\cB/N\) is sufficiently small.
The fixed marginal density bound permits deletion of these
exceptions.  Normalize the retained law, still denoted \(\rho\).
Its marginal bound \(\rho_i\le L\mu\) has \(L\) bounded in terms
of the early parameters and \(p_0\).

Project both modes of each component modulo \(T_e^\pm\).
Because \(\Delta\) is covered, the two endpoint tensors have the
same quotient.  The projections are injective on the individual
primal spans, so they preserve each tensor rank.  Choose a shortest
outer-product factorization of the common quotient tensor and
lift each factor to the two endpoint primal spans.  Componentwise,
the difference has the form
\begin{equation}
 \Delta=\sum_\alpha
   \bigl(p_\alpha b_\alpha\trans+
         a_\alpha q_\alpha\trans+
         a_\alpha b_\alpha\trans\bigr),
 \quad
 a_\alpha\in T^+,\quad b_\alpha\in T^- .
 \label{eq:cover-lifting}
\end{equation}
The first endpoint is the reference here; its factors are
\(p_\alpha,q_\alpha\), and those at the second endpoint are
\(p_\alpha+a_\alpha,q_\alpha+b_\alpha\).
The factor lists at either endpoint are independent in their
respective modes.

Choose bases for the spans of the offsets \(a_\alpha,b_\alpha\)
in each component.  Their combined ordered list is denoted
\(\omega_O\), and its size is \(t_O\), with
\begin{equation}
 1\le t_O\le2r.
 \label{eq:offset-rank}
\end{equation}
The lower bound follows from \(\Delta\ne0\).
Collect the companions of these offset basis vectors in a tuple
\(Z_O\): minus companions pair with plus offsets and vice versa.
The companions are independent within each component and mode.
For example, writing \(b_\alpha=\sum_kc_{k\alpha}\omega_k^-\)
gives a full-row-rank matrix \((c_{k\alpha})\), so the vectors
\(\sum_\alpha c_{k\alpha}p_\alpha\) are independent.
Using the second endpoint as reference changes each companion
by a known combination of offsets, and preserves independence.

An offset record consists of these basis vectors, their component
and sign assignments, and the bounded coefficient arrays in
\Cref{eq:cover-lifting}.  Its logarithmic count is at most
\[
 2r^2s+O_{r,g}(1)<.002N
\]
for sufficiently large \(n\).  Choosing a nominal representation
of a companion tuple in one endpoint's primal frame costs
\(2^{O(r\dim\cB)}=2^{O(n)}\) possibilities.  This last cost has
an arbitrarily small ratio exponent after choosing \(M_0\).
Fix deterministic conventions for all these records and
factorizations; their choices can be arbitrary functions of the
marked unit.

For a linear tensor functional \(u\), let \(\Phi_u(\omega)\)
be the tuple obtained by contracting \(u\) with the offset basis
vectors.  In particular, on a plus offset \(a\) its corresponding
plus output is \(X(Y\trans a)\); on a minus offset \(b\) the
minus output is \(Y(X\trans b)\).  At fixed two offset records
the phase in the theorem is, up to a sum of two unary phases,
\begin{equation}
 Z_A\cdot\Phi_{u_{B,S}}(\omega_A)
 +Z_B\cdot\Phi_{u_{A,R}}(\omega_B).
 \label{eq:offset-walsh}
\end{equation}
The unary phases are evaluations of the last term of
\Cref{eq:cover-lifting}; fixed offset records make them separate
functions of \(A\) and \(B\).

\medskip\noindent
\emph{Singleton phase estimates.}
All record laws in the following calculation are unnormalized
subprobabilities of \(\rho\).  For any record, the companion
tuple alone has point masses at most
\begin{equation}
 p_{\max}(Z_O)\le2^{-t_ON+.01N}.
 \label{eq:companion-pmax}
\end{equation}
To obtain this, enumerate its nominal primal representations,
apply the individual frame-image estimate and the marginal
bound \(L\mu\), and choose \(M_0\) large enough to absorb the
coefficient count and lost primal codimensions.

Suppose \(S=\{j\}\), and use endpoint \(j\) as the reference
for the \(B\)-companions.  For fixed \(\omega_A\), discard
pairs on which the coefficient lists
\(Y_j\trans a\) or \(X_j\trans b\), for its independent offsets,
fail full rank in any component.  An individual channel frame
is uniformly located by invariance.  Comparing it to iid
columns, the transpose-rank failure on \(d\) independent
ambient vectors is at most \(2^{d-h+1}\) for large \(n\).
Consequently the discarded pair mass is at most
\[
 L\,O(r)\,2^{2r-h}<10^{-4},
\]
on making \(h\) sufficiently large using early parameters only.

On the remaining pairs, the tuple
\((Z_B,\Phi_{u_{B,j}}(\omega_A))\) has point masses at most
\begin{equation}
 2^{-(t_B+t_A)N+.02N}.
 \label{eq:singleton-joint-pmax}
\end{equation}
Here is a direct conditional count.  Enumerate a nominal
representation of \(Z_B\), first test its values, and condition
on the primal frames at endpoint \(j\).  The two channel maps
are independent uniform-column draws in the appropriate
annihilators, apart from a bounded injectivity conditioning.
Freeze the values of \(Y\trans a\) and \(X\trans b\).  Their
number of possibilities is bounded in \(n\), and their lists
have full rank on the retained event.  Ignore the equations
that originally defined these coefficient values.  Prescribing
their images under \(X\) and \(Y\) then costs at least
\((N-\dim\cB)t_A\) bits, up to a bounded factor.
Together with the primal companion test this gives
\Cref{eq:singleton-joint-pmax}.  The ignored equations only
enlarge the event, so no independence between a coefficient
list and the map defining it has been assumed.

Apply \Cref{lem:walsh} to \Cref{eq:offset-walsh}, on
\((t_A+t_B)N\) bits, with the two tuples ordered as
\[
 \bigl(Z_A,\Phi_{u_{A,R}}(\omega_B)\bigr),
 \qquad
 \bigl(\Phi_{u_{B,j}}(\omega_A),Z_B\bigr).
\]
On the \(A\)-side the point-mass bound for the entire tuple is no
larger than
\Cref{eq:companion-pmax}; on the \(B\)-side use
\Cref{eq:singleton-joint-pmax}.  Separate unary phases have
absolute value one.  The contribution for two fixed records is
at most
\[
 2^{(t_A+t_B)N/2}
 \left(2^{-t_AN+.01N}
       2^{-(t_A+t_B)N+.02N}\right)^{1/2}
 =2^{(-t_A/2+.015)N}.
 \]
There are at most \(2^{.004N}\) record pairs.  By
\Cref{eq:offset-rank}, their total is at most \(2^{-.481N}\).
Thus every singleton phase mean is bounded by
\(10^{-4}+2^{-.481N}\).  The same proof applies when
\(|R|=1\).

\medskip\noindent
\emph{A tiny cover or mixing of the full endpoint sum.}
Put \(\delta=.0005\).  If a cover of total dimension at most
\(d_0\) contains the whole offset list with probability at least
\(\delta/4\), restrict to those units.  Delete units whose
channel transpose at either endpoint is not injective on the
fixed bases of that cover.  The deletion has arbitrarily small
fixed relative mass when \(h\) is large; the required lower
bound for \(h\) uses only \(L,\delta,d_0\).  On the remaining law
every offset list passes the singleton rank test, so the
preceding calculation has no pair discard and gives
\(2^{-\Omega(N)}\) for all singleton means.  Every \(\Delta\)
is covered by this fixed tiny cover.  This is the second
alternative.

Assume instead that no such cover captures mass \(\delta/4\).
Fix an \(A\)-offset list \(v=\omega_A\), and use endpoint \(1\)
as reference throughout the next calculation.  Call \(B\)
frequent for this list if its offset record \(m\), its companion
value \(z\), and its output \(\phi=\Phi_{u_{B,\{1,2\}}}(v)\)
have joint point probability greater than
\(2^{-(t_B+.5)N}\).  For fixed \(m,z\), define the deterministic
list
\begin{equation}
 {\cal F}_v(m,z)=
 \left\{\phi:
  \rho\bigl(m_B=m,Z_B=z,
       \Phi_{u_{B,\{1,2\}}}(v)=\phi\bigr)
       >2^{-(t_B+.5)N}\right\}.
 \label{eq:frequent-list}
\end{equation}
By \Cref{eq:companion-pmax},
\(\abs{{\cal F}_v(m,z)}\le2^{.51N}\).

We claim that the frequent-pair mass is at most \(\delta\).
Otherwise a set of \(B\)'s of mass at least \(\delta/2\) has
frequent \(A\)-mass at least \(\delta/2\).  Draw \(L_0\)
independent \(A\)-probes.  At each stage the spans of all past
offsets have total dimension at most \(2rL_0<d_0\).  The mass
of a probe whose whole offset list lies in those spans is
less than \(\delta/4\).  Thus each stage has conditional
probability at least \(\delta/4\) of being frequent and having
a fresh offset vector.  Select one fresh vector by a fixed
rule using only the probe lists.  Averaging fixes a probe
collection with independent selected witnesses in each
component and mode, and with successful \(B\)-mass at least
\begin{equation}
 q_*=(\delta/2)(\delta/4)^{L_0}>0.
 \label{eq:fresh-probe-mass}
\end{equation}
This lower bound is independent of selector and channel sizes.
Require
\begin{equation}
 L\,O(L_0)\,2^{L_0-h}<q_*/2.
 \label{eq:early-h-probes}
\end{equation}
The successful \(B\)'s can then be trimmed so that the channel
transpose at endpoint \(1\) has full rank on all these fixed
witness vectors, while leaving mass at least \(q_*/2\).

Now estimate this event under the raw law \(\mu^2\).
Condition on all of endpoint \(2\) and on the primal frames
at endpoint \(1\).  Enumerate the offset record \(m\), and
enumerate a nominal primal tuple representing \(Z_B\) at
endpoint \(1\).  There are at most \(2^{.002N+O(n)}\) such
choices.  For each choice the ambient value \(z\), all the
lists in \Cref{eq:frequent-list}, and their endpoint-\(2\)
translations are fixed.  The event defined by the adaptive
marks is contained in the union of these raw events.  This
uses no conditional density assertion about the marked law,
and does not enumerate \(2^{t_BN}\) arbitrary ambient
companion values.

Project each of the \(L_0\) lists onto the coordinate of its
chosen fresh witness.  Their total product size is at most
\(2^{.51L_0N}\).  Freeze the bounded lists of channel
coefficient values on the witnesses.  On their full-rank
event, prescribing the corresponding endpoint-\(1\) output
tuple costs at least \(.99L_0N\) bits, by the same annihilator
count as in \Cref{eq:singleton-joint-pmax}.  The upper
probability under \(\rho\), after paying its joint density,
is therefore at most
\[
 2^{(D+.01)N+.002N+O(n)-.48L_0N+O(1)}.
\]
Choose \(M_0\) to make the \(O(n)/N\) term small.
Because \(L_0\ge100(D+10)\), this is exponentially small,
contradicting the mass \(q_*/2\).  The claim follows.

Discard the frequent events in a full-sum phase test.
For fixed records, this is a restriction on \(B\) alone once
the \(A\)-offset record is fixed.  The \(B\)-tuple point bound
is now \(2^{-(t_B+.5)N}\); the \(A\)-tuple bound is
\(2^{-(t_A-.01)N}\).  Walsh on \((t_A+t_B)N\) bits gives
\(2^{-.245N}\), before the \(2^{.004N}\) record count.
Thus every phase with \(S=\{1,2\}\) has magnitude at most
\(\delta+2^{-\Omega(N)}\).  Interchanging the units covers
the cases with \(R=\{1,2\}\).  Together with the singleton
estimates, all nontrivial means are less than \(.005\).

All the positive-mass restrictions made above have lower
bounds in terms of \(p_0,\delta\) and early tolerances.
The inequalities for \(h\) were imposed after those quantities
were fixed.  The \(h\)-dependent bounded factors in the raw
counts only affect how large \(n\) must be.  This completes
both alternatives and the parameter-order assertion.
\end{proof}

The remaining phase conditions may hold on only inverse-polynomial
pair mass.  The next lemma makes injection failure small relative
to that mass.  It estimates the accepting pair event under independent
draws, without conditioning the two units jointly on acceptance.

\begin{lemma}[Injection on accepting pairs]
\label{lem:injection}
Let a prepared marked-unit law have final pins of total dimension
at most \(K\), and mixed marginals at most \(M2^{o(N)}\mu\).
Let \({\cal A}\) be a pair event for independent draws, of probability
at least \(N^{-c_0}\) for some fixed \(c_0\).
For each fixed draw suppose its accepting opposite units form a
member of a fixed family of cardinality at most \(2^{C_sN}\).
Assume \(C_s\) is fixed before \(h\).
For every fixed \(\eps>0\), by choosing \(h\) sufficiently large
in terms of \(K,C_s,\eps,g\), the proportion of pairs in
\({\cal A}\) failing the small-table injection conditions is at
most \(\eps\), for all sufficiently large \(n\).
For the unconditional pair event the accepting family has one
member.
\end{lemma}

\begin{proof}
It suffices to prove the assertion for a fixed component, endpoint
pair, and sign, with a sufficiently small failure tolerance
\(\eta>0\), and then sum over the at most \(16|\cE|\) tests.
For example, let \(V_A\) be the image of the individual pinned
primal space \(D_{A,e}^+\cap\cB_i\), and let \(H_B\subseteq\F^h\)
be the protected coefficient subspace at the tested endpoint
of \(B\).  Both dimensions are at most \(K\).  Injection fails
exactly when some nonzero \(v\in V_A\) has
\(Y_B\trans v\in H_B\).

Fix an accepting-set member \(S\) of \(A\)-mass at least
\(2^{-.01N}\), and put
\(\ell=\lfloor N/(10(K+1))\rfloor\).
Under \(A\) conditioned on \(S\), a new individual primal
frame span meets any previously fixed space of dimension at
most \(K\ell\) with probability at most
\[
 M2^{o(N)+.01N}
 2^{\dim\cB+K\ell-N+1}=2^{-\Omega(N)}.
\]
This bound includes every possible nominal pin choice, since
it tests the entire individual primal frame span.
If a fixed \(B\) has failure probability greater than \(\eta\)
on \(S\), successive draws \(A_1,\ldots,A_\ell\) from \(S\)
can each witness failure while their entire spaces \(V_{A_j}\)
are mutually independent.  At every step the conditional
probability is at least \(c=\eta/2\), for large \(n\).
The probability of the complete test is therefore at least
\(c^\ell\).

Conversely fix any such independent spaces.  There are at most
\(2^{K\ell}\) ways of choosing a nonzero witness in each.
The chosen witnesses are independent.  For iid \(Y_B\) columns,
their transpose images are independent uniform elements of
\(\F^h\).  If \(H_B\) were fixed, their membership probability
would be at most \(2^{-(h-K)\ell}\).
There are at most \((K+1)2^{Kh}\) protected coefficient spaces
of dimension at most \(K\), a constant in \(n\).
Conditioning \(Y_B\) to have full rank costs a bounded factor.
Thus, even allowing adaptive \(H_B\), the raw probability of
failure throughout is at most
\[
 C_{K,h}2^{-(h-2K)\ell}.
\]
The \(B\)-marginal density multiplies this by \(M2^{o(N)}\).
Fubini and the preceding lower test bound imply
\[
 \PP_B\{\PP_{A\mid S}(\text{failure}\mid B)>\eta\}
 \le C_{K,h}M
 2^{-(h-2K-\log_2(1/c))\ell+o(N)}.
 \]
Choose, for example,
\[
 h>2K+\log_2(1/c)+10(K+1)(C_s+2).
 \]
Taking a union over the accepting-set family leaves an
exponentially small exceptional \(B\)-mass.
Accepting sets of mass less than \(2^{-.01N}\) contribute
at most \(2^{-.01N}\) to the pair event.  On the remaining
nonexceptional \(B\)'s, failures have mass at most \(\eta\)
times the accepting mass.  Thus this one test fails on at
most \(\eta\PP({\cal A})+2^{-\Omega(N)}\).
Choose \(\eta<\eps/(32|\cE|)\), sum all tests in both roles,
and use \(\PP({\cal A})\ge N^{-c_0}\).
The exponential errors are negligible relative to this
probability, which proves the result.
\end{proof}

\begin{lemma}[Consequences of a tiny cover]
\label{lem:tiny-cover}
Suppose the second alternative of \Cref{thm:phase-alternative}
holds, and write
\[
 F(A,B)=u_{B,\{1,2\}}(\Delta_A).
\]
Then the following assertions hold.
\begin{enumerate}[label=\textup{(\alph*)}]
\item The matrix \(F\) on the finite support has rank at most
\(d_0N\).  Its alternating part \(F+F\trans\) has rank at most
\(2d_0N\).  If
\(\PP\{F(A,B)+F(B,A)=0\}\to0\), there is a restriction of
mass \(2^{-o(N)}\) on which that alternating part is
identically zero.
\item If \(F(A,A)=0\) throughout, then for either fixed
\(c\in\F\),
\begin{equation}
 \PP\{u_{B,1}(\Delta_A)=u_{B,2}(\Delta_A)
       =u_{A,1}(\Delta_B)=u_{A,2}(\Delta_B)=c\}
       \ge c_1N^{-4}
 \label{eq:tiny-four-phases}
\end{equation}
for a positive fixed \(c_1\) and all large \(n\).
\item In part \textup{(b)}, one can apply the second peeling
so that the event in \Cref{eq:tiny-four-phases} is determined
by the two leaves.  At most \(4K_1+2d_0<u_0\) additional
image directions are needed.  The same \(\Omega(N^{-4})\)
leaf-pair mass survives the exponentially small discards.
On admissible small tables, the contractions on the old
\(\lambda_i\)'s then equal the actual contractions.
\end{enumerate}
\end{lemma}

\begin{proof}
Write \((T_{0,e}^+,T_{0,e}^-)_{e\in\cE}\) for the fixed cover.
The cover sum space has dimension at most
\(\sum_e(\dim T_{0,e}^++\dim T_{0,e}^-)N\le d_0N\).
Since \(F\) pairs a vector \(\Delta_A\) in this space with a
functional determined by \(B\), the stated rank bounds follow.

For part (a), let \(\eps_N\) be the zero probability of the
alternating kernel \(H=F+F\trans\), and retain those rows whose
zero probability is at most \(\sqrt{\eps_N}\).  Their mass is
at least \(1-\sqrt{\eps_N}\).  Under the normalized retained
law, each such row is zero with probability at most
\(\sqrt{\eps_N}/(1-\sqrt{\eps_N})\).
Choose a row basis for the restricted kernel, of length at
most \(2d_0N\).  The entropy of its list of evaluations on a
random retained unit is bounded using binary entropy
\[
 H_2(t)=-t\log_2t-(1-t)\log_2(1-t),\qquad 0\log_2 0=0.
\]
It is at most
\[
 2d_0N\,
 H_2\left(\frac{\sqrt{\eps_N}}{1-\sqrt{\eps_N}}\right)=o(N).
 \]
A most likely list therefore has mass \(2^{-o(N)}\).
On its fiber, every row of the restricted kernel has the
same value on every two columns, because the chosen rows
span all rows.  Taking one column to be the row index itself
shows that this value is zero: \(H(A,A)=0\).  This proves (a).

For part (b), set \(R_N=1+d_0N\).  A pairwise incompatible
subset of the support, meaning that no distinct pair has
\(F(A,B)=F(B,A)=0\), gives
\[
 (1+F)\mathbin{\odot}(1+F)\trans=I
 \]
on that subset.  The entrywise-product rank is at most
\(\rank(1+F)^2\le R_N^2\), by expanding both factors as sums
of rank-one matrices.  Thus the subset has at most \(R_N^2\)
members.  Take \(R_N^2+1\) independent samples from the
weighted law.  A repeated sample is compatible because
\(F(A,A)=0\); otherwise the samples cannot all be pairwise
incompatible.  A union bound gives
\begin{equation}
 \PP\{F(A,B)=F(B,A)=0\}
 \ge \binom{R_N^2+1}{2}^{-1}.
 \label{eq:tiny-polynomial}
\end{equation}

Apply Fourier inversion to the four bits in
\Cref{eq:tiny-four-phases}.  The four characters with even
support in each endpoint pair contribute exactly
one quarter of the probability in \Cref{eq:tiny-polynomial}.
Each of the other twelve characters contains a singleton
subset and has magnitude \(2^{-\Omega(N)}\).
The coefficients of the even characters are one for either
common prescribed value \(c\).  Thus
\[
 \PP(\text{all four bits equal }c)
 =\frac14\PP\{F(A,B)=F(B,A)=0\}
   +O(2^{-\Omega(N)}),
 \]
which proves \Cref{eq:tiny-four-phases}.

Finally record the choices of the old \(\lambda_i\)'s in their
effective bases and pin the primal support images needed to
determine their tensors at both endpoints.  A bound of \(4K_1\)
directions suffices.  On each fixed plus cover basis vector
\(p\), and at each endpoint, store the bounded coefficient
list \(Y\trans p\) and pin the channel image \(X(Y\trans p)\).
For a minus cover basis vector \(q\), store \(X\trans q\)
and pin \(Y(X\trans q)\).  There are at most \(2d_0\) such
image directions over the two endpoints.  These data determine
each \(\Delta\) and both \(u\)-functionals on the whole cover
sum space, because
\[
 u(pq\trans)=\bigl(X(Y\trans p)\bigr)\cdot q
             =p\cdot\bigl(Y(X\trans q)\bigr).
 \]
They therefore determine the four phase bits for any pair
of leaves.

The total new rank is at most \(4K_1+2d_0<u_0\).
The scalar channel lists have bounded length in \(n\);
nominal support descriptions cost \(O(n)\) bits.
\Cref{lem:peeling} permits this refinement and the subsequent
second peeling, with total final rank at most \(K\), losing
only exponentially small total mass.  Removing unit mass
\(\varepsilon_N\) changes independent pair mass by at most
\(2\varepsilon_N\), which is negligible in
\Cref{eq:tiny-four-phases}.  The old support directions
are now pinned, so every admissible table evaluates their
tensors by actual frozen entries.  This proves (c).
\end{proof}

\begin{remark}
\label{rem:tiny-accepting-family}
For later applications of \Cref{lem:injection}, the four-phase
event in \Cref{eq:tiny-four-phases}, or a condition on
\(F(A,B)+F(B,A)\), has an accepting family of log-size
\(C_sN\), with \(C_s\le100(r+d_0+1)\).
Indeed the opposite unit can be described, for this event,
by its rank-at-most-\(r\) tensor \(\Delta\) and its two
functional restrictions to the cover sum space.
The former has at most \(2rN+O(1)\) bits, and the latter
at most \(2d_0N\) bits.  Enumerating all such descriptions
supplies a fixed family containing every accepting set.
The bound is independent of \(h\), of nominal leaf bases,
and of the selector dimensions.
\end{remark}
